\documentclass[11pt,a4paper]{article}

% Idioma y caracteres
\usepackage[utf8]{inputenc}
\usepackage[T1]{fontenc}
\usepackage[spanish]{babel}

% Diseño
\usepackage[margin=2.5cm]{geometry}
\usepackage{fancyhdr}
\usepackage{tabularx}
\usepackage{longtable}
\usepackage[table]{xcolor}
\usepackage{hyperref}

\setlength{\headheight}{14pt}
\setlength{\parindent}{0pt}
\setlength{\parskip}{6pt}
\renewcommand{\arraystretch}{1.3}

% Encabezado y pie de página
\pagestyle{fancy}
\fancyhf{}
\fancyhead[L]{\small UTEQ · Ingeniería de Requerimientos}
\fancyhead[R]{\small FORM-TL-01}
\fancyfoot[L]{\small Angelo Espinoza}
\fancyfoot[R]{\small Página \thepage}

\begin{document}

\begin{center}
    {\large\textbf{
        UNIVERSIDAD TÉCNICA ESTATAL DE QUEVEDO
    }}

    Facultad de Ciencias de la Computación\\
    Carrera de Software

    \vspace{0.6cm}

    {\LARGE\textbf{Taller de tipos de requisitos}}

    \vspace{0.2cm}

    {\large FORM-TL-01}\\
    Actividad formativa
\end{center}

\vspace{0.4cm}

\section{Datos informativos}

\begin{tabularx}{\textwidth}{|l|X|}
\hline
\rowcolor{gray!15}
\textbf{Campo} & \textbf{Detalle} \\
\hline
Estudiante & Angelo Maximiliano Espinoza Pivaque \\
\hline
Asignatura & Ingeniería de Requerimientos (ISR-401) \\
\hline
Nivel y paralelo & Cuarto nivel, paralelo A \\
\hline
Docente & Gleiston Cicerón Guerrero Ulloa, PhD \\
\hline
Caso & Plataforma de gestión de citas médicas \\
\hline
\end{tabularx}

\section{Clasificación de los enunciados}

Los enunciados se clasifican en tres dimensiones:
tipo de requisito, producto o proceso, y sistema
o software. Para los requisitos de calidad se
incluye la característica y subcaracterística
de ISO/IEC 25010:2023.

\begin{longtable}{|
    p{0.08\textwidth}|
    p{0.82\textwidth}|}

\hline
\rowcolor{gray!15}
\textbf{Código} & \textbf{Clasificación y justificación} \\
\hline
\endfirsthead

\hline
\rowcolor{gray!15}
\textbf{Código} & \textbf{Clasificación y justificación} \\
\hline
\endhead

\textbf{E01} &
\textbf{Enunciado:}
El paciente podrá reservar una cita eligiendo
especialidad, médico, fecha y hora entre los
horarios disponibles.

\textbf{Tipo:} Funcional.

\textbf{ISO/IEC 25010:2023:}
No aplica como clasificación de calidad.

\textbf{Producto o proceso:} Producto.

\textbf{Sistema o software:} Software.

\textbf{Justificación:}
Define la función de reservar citas que debe
ofrecer la aplicación al paciente.
\\
\hline

\textbf{E02} &
\textbf{Enunciado:}
La plataforma enviará al paciente un recordatorio
por mensaje de texto 24 horas antes de su cita.

\textbf{Tipo:} Funcional.

\textbf{ISO/IEC 25010:2023:}
No aplica como clasificación de calidad.

\textbf{Producto o proceso:} Producto.

\textbf{Sistema o software:} Software.

\textbf{Justificación:}
Especifica el envío automático de un recordatorio
y el momento en que debe ejecutarse.
\\
\hline

\textbf{E03} &
\textbf{Enunciado:}
La aplicación debe ser rápida.

\textbf{Tipo:} Calidad.

\textbf{ISO/IEC 25010:2023:}
Eficiencia de desempeño
(\textit{performance efficiency});
comportamiento temporal
(\textit{time behaviour}).

\textbf{Producto o proceso:} Producto.

\textbf{Sistema o software:} Software.

\textbf{Justificación:}
Expresa una expectativa sobre el tiempo de respuesta
de la aplicación, aunque no establece un umbral
que permita verificarla.
\\
\hline

\textbf{E04} &
\textbf{Enunciado:}
El recepcionista podrá registrar la llegada del
paciente que se acerca a la ventanilla.

\textbf{Tipo:} Funcional.

\textbf{ISO/IEC 25010:2023:}
No aplica como clasificación de calidad.

\textbf{Producto o proceso:} Producto.

\textbf{Sistema o software:} Software.

\textbf{Justificación:}
Define la función de registrar en la aplicación
la llegada de un paciente.
\\
\hline

\textbf{E03} &
\textbf{Enunciado:}
La aplicación debe ser rápida.

\textbf{Tipo:} Calidad.

\textbf{ISO/IEC 25010:2023:}
Eficiencia de desempeño
(\textit{performance efficiency});
comportamiento temporal
(\textit{time behaviour}).

\textbf{Producto o proceso:} Producto.

\textbf{Sistema o software:} Software.

\textbf{Justificación:}
Expresa una expectativa sobre el tiempo de respuesta
de la aplicación, aunque no establece un umbral
que permita verificarla.
\\
\hline

\textbf{E04} &
\textbf{Enunciado:}
El recepcionista podrá registrar la llegada del
paciente que se acerca a la ventanilla.

\textbf{Tipo:} Funcional.

\textbf{ISO/IEC 25010:2023:}
No aplica como clasificación de calidad.

\textbf{Producto o proceso:} Producto.

\textbf{Sistema o software:} Software.

\textbf{Justificación:}
Define la función de registrar en la aplicación
la llegada de un paciente.
\\
\hline


\textbf{E05} &
\textbf{Enunciado:}
La plataforma se desarrollará en Java 21 con
base de datos PostgreSQL, tecnologías que
domina el área de TI.

\textbf{Tipo:} Restricción tecnológica.

\textbf{ISO/IEC 25010:2023:}
No aplica como clasificación de calidad.

\textbf{Producto o proceso:} Producto.

\textbf{Sistema o software:} Software.

\textbf{Justificación:}
Limita las tecnologías de implementación del
software al exigir Java 21 y PostgreSQL.
\\
\hline

\textbf{E06} &
\textbf{Enunciado:}
El médico podrá consultar la lista de pacientes
del día y marcar cada consulta como atendida.

\textbf{Tipo:} Funcional.

\textbf{ISO/IEC 25010:2023:}
No aplica como clasificación de calidad.

\textbf{Producto o proceso:} Producto.

\textbf{Sistema o software:} Software.

\textbf{Justificación:}
Define las funciones de consultar la agenda
diaria y registrar las consultas atendidas.
\\
\hline

\textbf{E07} &
\textbf{Enunciado:}
Los datos de salud de los pacientes deben
estar seguros.

\textbf{Tipo:} Calidad.

\textbf{ISO/IEC 25010:2023:}
Seguridad. Se proponen confidencialidad e
integridad como subcaracterísticas relevantes;
el enunciado original no precisa cuáles exige.

\textbf{Producto o proceso:} Producto.

\textbf{Sistema o software:} Sistema.

\textbf{Justificación:}
Expresa una expectativa de protección de los
datos en el sistema completo, pero no define
controles ni criterios verificables.
\\
\hline

\textbf{E08} &
\textbf{Enunciado:}
La aplicación móvil debe funcionar en teléfonos
con Android 10 o superior.

\textbf{Tipo:} Restricción de plataforma.

\textbf{ISO/IEC 25010:2023:}
No aplica como clasificación de calidad
en el criterio adoptado.

\textbf{Producto o proceso:} Producto.

\textbf{Sistema o software:} Software.

\textbf{Justificación:}
Impone el sistema operativo y la versión mínima
que debe admitir la aplicación móvil.
\\
\hline

\textbf{E09} &
\textbf{Enunciado:}
Una cita cancelada con menos de cuatro horas
de anticipación se registrará como inasistencia.

\textbf{Tipo:} Funcional, derivado de una
regla de negocio.

\textbf{ISO/IEC 25010:2023:}
No aplica como clasificación de calidad.

\textbf{Producto o proceso:} Producto.

\textbf{Sistema o software:} Software.

\textbf{Justificación:}
Define el comportamiento del software ante una
cancelación tardía: registrar una inasistencia.
\\
\hline

\textbf{E10} &
\textbf{Enunciado:}
La plataforma debe ser fácil de usar para
los pacientes adultos mayores.

\textbf{Tipo:} Calidad.

\textbf{ISO/IEC 25010:2023:}
Capacidad de interacción
(\textit{interaction capability});
operabilidad (\textit{operability}) e
inclusividad (\textit{inclusivity}).

\textbf{Producto o proceso:} Producto.

\textbf{Sistema o software:} Sistema.

\textbf{Justificación:}
Exige facilidad de operación para adultos
mayores en las interfaces de la plataforma,
incluido el quiosco, pero no define criterios
medibles de aceptación.
\\
\hline

\textbf{E11} &
\textbf{Enunciado:}
El proyecto debe entregarse en un plazo
máximo de seis meses.

\textbf{Tipo:} Restricción temporal.

\textbf{ISO/IEC 25010:2023:}
No aplica; establece un plazo del proyecto,
no una propiedad de calidad del producto.

\textbf{Producto o proceso:} Proceso.

\textbf{Sistema o software:}
Sistema, por referirse a la entrega del
proyecto completo.

\textbf{Justificación:}
Limita el tiempo disponible para desarrollar
y entregar la solución completa.
\\
\hline

\textbf{E12} &
\textbf{Enunciado:}
La sesión de un usuario se cerrará
automáticamente tras 10 minutos sin actividad.

\textbf{Tipo:}
Funcional, con propósito de seguridad.

\textbf{ISO/IEC 25010:2023:}
No aplica como clasificación de calidad;
especifica un comportamiento concreto
que contribuye a la seguridad.

\textbf{Producto o proceso:} Producto.

\textbf{Sistema o software:} Software.

\textbf{Justificación:}
Define una acción automática del software:
cerrar la sesión al cumplirse 10 minutos
de inactividad.
\\
\hline

\textbf{E13} &
\textbf{Enunciado:}
El quiosco de la sala de espera leerá el número
de cédula ingresado en su pantalla táctil e
imprimirá un ticket con el número de turno.

\textbf{Tipo:} Funcional.

\textbf{ISO/IEC 25010:2023:}
No aplica como clasificación de calidad.

\textbf{Producto o proceso:} Producto.

\textbf{Sistema o software:} Sistema.

\textbf{Justificación:}
Define funciones que requieren la interacción
del software con la pantalla táctil y la
impresora térmica del quiosco.
\\
\hline

\textbf{E14} &
\textbf{Enunciado:}
La plataforma debe estar siempre disponible.

\textbf{Tipo:} Calidad.

\textbf{ISO/IEC 25010:2023:}
Fiabilidad (\textit{reliability});
disponibilidad (\textit{availability}).

\textbf{Producto o proceso:} Producto.

\textbf{Sistema o software:} Sistema.

\textbf{Justificación:}
Exige que el servicio esté operativo y accesible,
lo que depende del software y la infraestructura;
no define un periodo ni un porcentaje de
disponibilidad para verificar su cumplimiento.
\\
\hline
\textbf{E15} &
\textbf{Enunciado:}
El director médico podrá generar un reporte
mensual de inasistencias por especialidad.

\textbf{Tipo:} Funcional.

\textbf{ISO/IEC 25010:2023:}
No aplica como clasificación de calidad.

\textbf{Producto o proceso:} Producto.

\textbf{Sistema o software:} Software.

\textbf{Justificación:}
Define la función de generar reportes de
inasistencias agrupadas por especialidad
para un periodo mensual.
\\
\hline

\textbf{E16} &
\textbf{Enunciado:}
Cada incremento se validará con el director
médico al final de cada iteración de dos semanas.

\textbf{Tipo:} Restricción del proceso.

\textbf{ISO/IEC 25010:2023:}
No aplica; establece una condición del proceso
de desarrollo, no de calidad del producto.

\textbf{Producto o proceso:} Proceso.

\textbf{Sistema o software:}
Sistema, por abarcar los incrementos de la
solución completa.

\textbf{Justificación:}
Impone quién participa en la validación de
cada incremento y la periodicidad de
esa actividad.
\\
\hline

\textbf{E17} &
\textbf{Enunciado:}
El tratamiento de los datos personales de
los pacientes debe cumplir la Ley Orgánica
de Protección de Datos Personales del Ecuador.

\textbf{Tipo:} Restricción legal.

\textbf{ISO/IEC 25010:2023:}
No aplica como clasificación de calidad;
el cumplimiento legal no se reduce a
una característica del modelo.

\textbf{Producto o proceso:} Producto,
con implicaciones para el proceso.

\textbf{Sistema o software:} Sistema.

\textbf{Justificación:}
Impone un marco legal al tratamiento de
datos personales en la solución completa,
incluidos sus procedimientos de operación.
\\
\hline

\textbf{E18} &
\textbf{Enunciado:}
La plataforma debe integrarse sin problemas
con el sistema de facturación actual.

\textbf{Tipo:} Calidad.

\textbf{ISO/IEC 25010:2023:}
Compatibilidad (\textit{compatibility});
interoperabilidad (\textit{interoperability}).

\textbf{Producto o proceso:} Producto.

\textbf{Sistema o software:} Sistema.

\textbf{Justificación:}
Exige el intercambio y uso de información
entre la plataforma y el sistema de facturación,
pero no precisa qué significa integrarse
sin problemas.
\\
\hline

\textbf{E19} &
\textbf{Enunciado:}
Si el servidor principal falla, la plataforma
debe restablecer el servicio en menos de
30 minutos sin perder ninguna cita confirmada.

\textbf{Tipo:} Calidad.

\textbf{ISO/IEC 25010:2023:}
Fiabilidad (\textit{reliability});
capacidad de recuperación
(\textit{recoverability}).

\textbf{Producto o proceso:} Producto.

\textbf{Sistema o software:} Sistema.

\textbf{Justificación:}
Establece condiciones de recuperación del
servicio y de los datos ante un fallo del
servidor: menos de 30 minutos y ninguna
cita confirmada perdida.
\\
\hline

\textbf{E20} &
\textbf{Enunciado:}
Todo cambio en el código fuente será revisado
por otro integrante del equipo antes de
integrarse a la rama principal.

\textbf{Tipo:} Restricción del proceso.

\textbf{ISO/IEC 25010:2023:}
No aplica; exige una práctica de desarrollo,
no una propiedad de calidad del producto.

\textbf{Producto o proceso:} Proceso.

\textbf{Sistema o software:} Software.

\textbf{Justificación:}
Impone la revisión del código por otro
integrante como condición previa a su
integración en la rama principal.
\\
\hline
\end{longtable}

\section{Cuantificación de requisitos de calidad}



\section{Discrepancias y resolución}



\section{Referencias}



\section{Declaración de uso de inteligencia artificial}



\end{document}